\documentclass[11pt]{article}
\usepackage[margin=1in]{geometry}
\usepackage{amsmath,amssymb,amsthm}
\usepackage[T1]{fontenc}
\usepackage{lmodern}
\usepackage[hidelinks]{hyperref}
\newtheorem{theorem}{Theorem}
\newtheorem{lemma}{Lemma}
\newcommand{\R}{\mathbb R}
\newcommand{\Cent}{\mathcal C}
\title{The vertex centroid of focal antipedals\\of even elliptic billiards}
\author{Alec Kriebel\\\small\href{https://orcid.org/0009-0001-9320-500X}{ORCID: 0009-0001-9320-500X}}
\date{October 7, 2026}
\begin{document}
\maketitle
\begin{abstract}
For a periodic billiard in an ellipse with a strictly nested confocal
elliptical caustic and even least period, we prove that the unweighted
vertex centroids of the original antipedals about the center and the two
foci are constant throughout its Poncelet family. The central centroid
is zero; an explicit focal coefficient depends only on the axes,
caustic parameter, period and perimeter. This answers invariant k405
as posed without proof by Reznik, Garcia and Koiller. The proof uses
opposite edges, an elementary finite circle action, and perimeter
stationarity. Ferudun's subsequent deposited note states the same
theorem and coefficient and uses the same pair-and-moment mechanism;
we explicitly credit that overlap and make no absolute priority or
independent-discovery claim.
\end{abstract}

\section{Statement, source and conventions}
Let $a>b>0$, $c=\sqrt{a^2-b^2}$, and consider
\[
 E:\ \frac{x^2}{a^2}+\frac{y^2}{b^2}=1,
 \qquad E_\lambda:\ \frac{x^2}{a^2-\lambda}
                  +\frac{y^2}{b^2-\lambda}=1,
 \qquad 0<\lambda<b^2.
\]
Write $O=(0,0)$ and $F_\sigma=(\sigma c,0)$, $\sigma=\pm1$.
A periodic billiard tangent to $E_\lambda$ extends to the smooth
Poncelet family with the same caustic; we use this classical closure
theorem in its usual confocal-ellipse form \cite{Stachel}.
For successive billiard vertices $A,B$ and a pole $M$, define
$Q_M(A,B)$ as the intersection of the full lines through $A$ and $B$
perpendicular to $A-M$ and $B-M$, respectively. For an orbit of least
period $N$, its \emph{antipedal vertex centroid} is
\begin{equation}\label{defcent}
 \Cent_M=\frac1N\sum_{i=0}^{N-1}Q_M(P_i,P_{i+1}),\qquad P_N=P_0.
\end{equation}
This is a vertex mean, not an area centroid or a mean of pedal feet.
The antipedal belongs to the original orbit, not to its outer tangent
polygon. All three poles lie strictly inside $E_\lambda$, so its tangent
chords do not pass through them. The intersections in \eqref{defcent}
are therefore finite and unique.

\begin{samepage}
\begin{theorem}\label{main}
Suppose the oriented billiard has even least period $N$ (hence $N\ge4$)
and caustic $E_\lambda$. Let $L$ be the perimeter of its Poncelet family,
and put
\[
 K=a^2b^2-\lambda c^2,\qquad
 H=\frac{a^2}{c^2}\left(1-\frac{bL}{2a\sqrt\lambda\,N}\right).
\]
Then throughout that family,
\begin{equation}\label{answer}
 \Cent_O=(0,0),\qquad
 \Cent_{F_\sigma}=
 \left(\sigma c\left[-1+\frac{KH}{a^2(b^2-\lambda)}\right],0\right).
\end{equation}
Both simple and primitive star orbits, and either traversal orientation,
are allowed.
\end{theorem}
\end{samepage}

The source is k405 in Table~5 of \cite{RGK21}, p.~348, and of
\cite{RGK20}, printed p.~7. It specifies even $N$ and poles $O,F_1,F_2$
in the confocal-ellipse setting, with proof and value left as questions.
The source does not formally define least period. Its use of opposite
vertices under central inversion in \cite[Section~3.7]{RGK20} supports
the primitive-period interpretation used explicitly here. Repeating an
even-period polygon leaves \eqref{defcent} unchanged; repeating an odd
polygon an even number of times does not put it under this theorem.
Degenerate two-bounce orbits, limiting degenerate caustics and hyperbolic
caustics are outside the statement. In a circle the foci coincide with
$O$ and the conclusion for even least period is zero by central symmetry;
the coefficient above, which divides by $c^2$, is not used for a circle.

The least-period restriction is substantive. For $a=2,b=1$, the
billiard triangle with vertices $(2,0)$ and
\[
 \left(\frac{2(1-\sqrt{13})}{3},
       \ \pm\frac{\sqrt{2\sqrt{13}-5}}{3}\right)
\]
has caustic $\lambda=(8\sqrt{13}-20)/9$ and origin antipedal centroid
$((35-5\sqrt{13})/24,0)\ne0$; the reflected triangle has the opposite
centroid. Repeating each twice gives six-entry lists with these same
centroids. Exact reflection, tangency and centroid checks for this
boundary example are supplied in the supplement.

\section{Opposite vertices and two chord moments}
\begin{lemma}\label{central}
For the hypotheses of Theorem~\ref{main},
$P_{i+N/2}=-P_i$.
\end{lemma}
\begin{proof}
Let $T$ be the Poncelet tangent map, choosing the directed chord with
the inner ellipse on its left. This is a smooth orientation-preserving
homeomorphism of the outer oval; the other tangent gives its inverse.
The Poncelet family premise gives $T^N=\mathrm{id}$, of order exactly
$N$. Central inversion $S(P)=-P$ preserves the chosen side of a directed
chord, so $S$ commutes with $T$.

The finite group generated by $S$ and $T$ preserves the average of the
pushforwards of arc-length probability measure over this group. This
measure is nonatomic and positive on every nonempty arc. Its cumulative
coordinate conjugates every group element to a rotation: an oriented
measure-preserving map preserves all arc lengths in this coordinate.
A rotation group has a unique element of order two. Both $S$ and
$T^{N/2}$ have order two, so they agree.
\end{proof}
Opposite antipedal intersections for the pole $O$ are $Q$ and $-Q$.
Lemma~\ref{central} already proves $\Cent_O=0$.

Choose the orientation with the caustic to the left. Since $O$ is inside
the caustic, $\det(A,B)>0$ for each directed edge. Lift its eccentric
endpoint angles as $s-\delta,s+\delta$, with $0<\delta<\pi/2$, and write
\[
 C=\cos s,\quad S=\sin s,\quad U=\cos\delta,\quad V=\sin\delta,
 \quad D^2=\frac{C^2}{a^2}+\frac{S^2}{b^2}.
\]
Then
\begin{align*}
 A&=(a(CU+SV),\ b(SU-CV)),\\
 B&=(a(CU-SV),\ b(SU+CV)).
\end{align*}
The chord line is $xC/a+yS/b=U$. Tangency to $E_\lambda$ gives
\begin{equation}\label{tangent}
 U^2=1-\lambda D^2,\qquad V^2=\lambda D^2.
\end{equation}
Its positive length is
\begin{equation}\label{length}
 \ell=2V\sqrt{a^2S^2+b^2C^2}=2ab\sqrt\lambda\,D^2.
\end{equation}

\begin{lemma}\label{moments}
The perimeter $L$ is constant on the Poncelet family, and its chord
coordinates satisfy
\begin{equation}\label{moments_eq}
 \frac1N\sum_i C_i^2=H,\qquad \sum_i S_iC_i=0.
\end{equation}
\end{lemma}
\begin{proof}
Under independent tangent displacements of the billiard vertices,
the first variation of perimeter is
\[
 dL=\sum_i (u_{i-1}-u_i)\mathbin{\cdot}dP_i,
 \qquad u_i=\frac{P_{i+1}-P_i}{|P_{i+1}-P_i|}.
\]
The reflection law makes $u_{i-1}-u_i$ normal to $E$, so $dL=0$.
This calculation does not require the polygon to be simple. Along the
smooth Poncelet family it proves that $L$ has zero derivative and is
constant on the connected circle of initial vertices.

Summing \eqref{length} and using
$D^2=b^{-2}-c^2C^2/(a^2b^2)$ yields the first identity in
\eqref{moments_eq}. For the second, use the allowed comparison
deformation $t_i\mapsto t_i+\epsilon$ of all eccentric vertex angles.
It keeps each $\delta_i$ fixed and changes $s_i$ to $s_i+\epsilon$.
The deformed chords need not preserve the caustic. Differentiate the
first expression in \eqref{length} \emph{before} imposing a caustic
constraint on the deformation. At the original billiard,
\[
 \left.\frac{d\ell_i}{d\epsilon}\right|_0
 =\frac{2V_i c^2 S_iC_i}{\sqrt{a^2S_i^2+b^2C_i^2}}
 =\frac{2\sqrt\lambda\,c^2}{ab}S_iC_i.
\]
Stationarity of $L$ gives the claimed sum, since $c,\lambda>0$.
\end{proof}

\section{The opposite-edge antipedal calculation}
For a focus $F=(h,0)$, $h=\pm c$, set
\[
 m=(aCU,bSU),\quad v=(-aSV,bCV),\quad d=m-F,
 \quad q=Q_F(A,B)-F.
\]
Thus $A=m-v$ and $B=m+v$. Adding and subtracting the equations of the
two antipedal lines gives
\begin{equation}\label{system}
 q\cdot d=R:=|d|^2+|v|^2,\qquad
 q\cdot v=T:=2d\cdot v.
\end{equation}
Its determinant is $bV(aU-hC)>0$. Indeed,
\begin{equation}\label{denominator}
 U^2-\frac{c^2C^2}{a^2}=(b^2-\lambda)D^2>0,
\end{equation}
so $aU>|hC|$.

For clarity, Cramer's rule gives
\[
 q_x=\frac{bCVR-bSUT}{bV(aU-hC)},\qquad
 q_y=\frac{(aCU-h)T+aSVR}{bV(aU-hC)}.
\]
For the opposite edge replace $(m,v)$ by $(-m,-v)$ in
\eqref{system}, while keeping $F$ fixed. Adding the two resulting
absolute intersections, with $h^2=c^2$, gives
\begin{align}\label{pair_intermediate}
 \bigl(Q_F(A,B)+Q_F(-A,-B)\bigr)_x
   &=\frac{2hb^2(C^2-U^2)}{b^2-\lambda},\\
 \bigl(Q_F(A,B)+Q_F(-A,-B)\bigr)_y
   &=\frac{-2CSah\,[c^2(C^2-U^2)-b^2]}
           {ba^2(U^2-c^2C^2/a^2)}.\nonumber
\end{align}
Here the first expression has used \eqref{denominator} and
$D^2=C^2/a^2+S^2/b^2$. The remaining simplifications are
\[
 c^2(C^2-U^2)-b^2=-KD^2,\qquad
 b^2(C^2-U^2)=-(b^2-\lambda)+\frac{K C^2}{a^2}.
\]
Consequently
\begin{equation}\label{pair}
 Q_F(A,B)+Q_F(-A,-B)
 =\left(-2h+\frac{2hK C^2}{a^2(b^2-\lambda)},
              \frac{2hK SC}{ab(b^2-\lambda)}\right).
\end{equation}
This local identity uses tangency but not closure or the reflection
law. By Lemma~\ref{central}, summing its left side over all edges counts
each antipedal intersection twice. Divide by $2N$ and apply
Lemma~\ref{moments}. This proves \eqref{answer}. Reversing traversal
only permutes the same full-line intersections, so it leaves the result
unchanged.

\section{Provenance, overlap and verification}
Central symmetry, perimeter stationarity and confocal billiard
conservation are classical; see, for example, \cite{BT,Stachel}.
The arguments needed here have been supplied in full. Earlier related
pedal-centroid and harmonic results in \cite{BT}, together with the
classical skew-hodograph correspondence discussed in \cite{GT}, yield
another reduction once the target-specific pair identity \eqref{pair}
is inserted. These ingredients alone are not a located earlier proof
of the focal antipedal theorem. Our bounded audit of primary literature
did not establish an earlier complete answer to k405; it does not
establish absolute priority.

A proof was publicly posted in PR140 on September 30, 2026
\cite{PR140}. Its archived original already contains the theorem and
the pair formula. Ferudun's \emph{Original Antipedal Centroids in
Elliptic Billiards} \cite{Ferudun} states the same scoped theorem and
the same explicit coefficient. Its specific Zenodo record was created
on October 1, 2026 at 23:58:07 UTC, with publication date October 2.
Its proof also combines opposite-edge antipedal pairs with chord
moments; a change to unit-normal coordinates gives the same identities.
We credit this exact overlap. The documented dates concern those
specific public records; they do not establish absolute priority,
independent discovery, or any inference of copying or collaboration.
The present note records a verified self-contained resolution of a
question posed without proof in the 2020/2021 source, rather than
claiming that the question remains unsolved at this note's publication.

The accompanying verification package contains explicit-guard rational
and symbolic checks of the local line systems, tangency and pair
identities, length differentiation, and finite-rotation parity controls.
Its finite computations support reproducibility; they do not prove
Poncelet closure for arbitrary parameters or replace the all-period
arguments above. Separate adversarial audits tested genuine elliptic
orbits, parameter boundaries, repeated odd polygons and deliberately
false identities. Audit results and limits are described in the
supplement.

\paragraph{AI assistance and review status.}
AI tools were used extensively in developing, drafting, reproducing
and adversarially auditing this work. Independent AI review passes
are not conventional human peer review. This preprint is unrefereed
and has not undergone conventional human peer review at publication.
The proof and supplementary programs are supplied for independent
checking.

\begin{thebibliography}{9}
\bibitem{RGK21} D. Reznik, R. Garcia and J. Koiller,
\emph{Fifty New Invariants of N-Periodics in the Elliptic Billiard},
Arnold Mathematical Journal \textbf{7} (2021), 341--355.
\href{https://doi.org/10.1007/s40598-021-00174-y}{doi:10.1007/s40598-021-00174-y}.
\bibitem{RGK20} D. Reznik, R. Garcia and J. Koiller,
\emph{Eighty New Invariants in the Elliptic Billiard},
arXiv:2004.12497v11 (2020).
\href{https://arxiv.org/abs/2004.12497v11}{arXiv:2004.12497v11}.
\bibitem{BT} M. Bialy and S. Tabachnikov,
\emph{Dan Reznik's identities and more}, European Journal of
Mathematics \textbf{8} (2022), 1341--1354; first published online 2020.
\href{https://doi.org/10.1007/s40879-020-00428-7}{doi:10.1007/s40879-020-00428-7}.
\bibitem{GT} E. Gutkin and S. Tabachnikov,
\emph{Billiards in Finsler and Minkowski geometries}, Journal of
Geometry and Physics \textbf{40} (2002), 277--301.
\href{https://doi.org/10.1016/S0393-0440(01)00039-0}{doi:10.1016/S0393-0440(01)00039-0}.
\bibitem{Stachel} H. Stachel,
\emph{The geometry of billiards in ellipses and their Poncelet grids},
Journal of Geometry \textbf{112} (2021), article 40.
\href{https://doi.org/10.1007/s00022-021-00606-2}{doi:10.1007/s00022-021-00606-2}.
\bibitem{PR140} A. Kriebel, \emph{5100023: prove the focal antipedal
centroid invariant k405}, public PR140 and proof archive, September 30, 2026.
\href{https://github.com/AlecKriebel/Math/pull/140}{github.com/AlecKriebel/Math/pull/140}.
\bibitem{Ferudun} A. Ferudun,
\emph{Original Antipedal Centroids in Elliptic Billiards}, version 1.0,
Zenodo (publication date October 2, 2026).
\href{https://doi.org/10.5281/zenodo.23092466}{doi:10.5281/zenodo.23092466}.
\end{thebibliography}
\end{document}
